% ============================================================
%  DISSERTATION OUTLINE  (REVISED)
%  Title : Simulating Random Graphs in Python
%  Module: MA-300: Undergraduate Project in Mathematics
%  Target: 7000-8000 words
% ============================================================

% ────────────────────────────────────────────────────────────
%  PREAMBLE
% ────────────────────────────────────────────────────────────
\documentclass[12pt, a4paper]{report}

% --- Encoding & fonts ---
\usepackage[T1]{fontenc}
\usepackage[utf8]{inputenc}
\usepackage{lmodern}

% --- Page layout ---
\usepackage[a4paper, top=2.5cm, bottom=2.5cm,
            left=3.0cm, right=2.5cm]{geometry}
\usepackage{setspace}
\onehalfspacing

% --- Mathematics ---
\usepackage{amsmath}
\usepackage{amssymb}
\usepackage{amsthm}
\usepackage{mathtools}

% --- Figures ---
\usepackage{graphicx}
\usepackage{float}
\usepackage{caption}
\captionsetup{font=small, labelfont=bf}

\graphicspath{{./}}

% --- Tables ---
\usepackage{booktabs}

% --- Code listings ---
\usepackage{listings}
\usepackage{xcolor}

\definecolor{codegray}{rgb}{0.5,0.5,0.5}
\definecolor{codeblue}{rgb}{0.15,0.35,0.65}
\definecolor{codegreen}{rgb}{0.1,0.5,0.1}
\definecolor{codebg}{rgb}{0.97,0.97,0.97}

\lstdefinestyle{python}{
  language        = Python,
  backgroundcolor = \color{codebg},
  basicstyle      = \ttfamily\footnotesize,
  keywordstyle    = \color{codeblue}\bfseries,
  commentstyle    = \color{codegray}\itshape,
  stringstyle     = \color{codegreen},
  numberstyle     = \tiny\color{codegray},
  numbers         = left,
  stepnumber      = 1,
  breaklines      = true,
  frame           = single,
  captionpos      = b
}

% --- Hyperlinks ---
\usepackage[colorlinks=true,
            linkcolor=blue!60!black,
            citecolor=blue!60!black,
            urlcolor=blue!60!black]{hyperref}

% --- Theorem environments ---
\theoremstyle{definition}
\newtheorem{definition}{Definition}[chapter]
\newtheorem{example}{Example}[chapter]

\theoremstyle{plain}
\newtheorem{theorem}{Theorem}[chapter]
\newtheorem{proposition}{Proposition}[chapter]

\theoremstyle{remark}
\newtheorem{remark}{Remark}[chapter]

% --- Shorthand math ---
\newcommand{\avg}[1]{\langle #1 \rangle}
\newcommand{\Prob}{\mathbb{P}}
\newcommand{\E}{\mathbb{E}}

% ────────────────────────────────────────────────────────────
%  DOCUMENT
% ────────────────────────────────────────────────────────────
\begin{document}

% ════════════════════════════════════════════════════════════
%  TITLE PAGE
% ════════════════════════════════════════════════════════════
\begin{titlepage}
  \centering
  \vspace*{3cm}
  {\LARGE\bfseries Simulating Random Graphs in Python:\par}
  \vspace{0.4cm}
  {\large\bfseries
    Erd\H{o}s--R\'{e}nyi, Watts--Strogatz, and
    Barab\'{a}si--Albert Models\par}
  \vspace{2.5cm}
  {\large Zachariah Attila Markusson\par}
  \vspace{0.5cm}
  {\normalsize Student Number: 2370278\par}
  \vspace{2cm}
  {\normalsize
    Submitted in partial fulfilment of the requirements\\
    for the module \textbf{MA-300: Undergraduate Project in
    Mathematics}\\[0.3cm]
    Department of Mathematics\\
    Swansea University\\[0.3cm]
    Academic Year 2025--2026\par}
  \vfill
  {\small Word count: XXXX}
\end{titlepage}

\chapter{Mathematical Foundations}
\label{chap:foundations}

In this chapter I will summarise the mathematical background needed
for the three random graph models studied later. I aim to explain basic
graph-theory definitions, the network metrics used, the spectral
viewpoint via the graph Laplacian, and the probabilistic concepts
underlying the Erd\H{o}s--R\'enyi model, the Watts--Strogatz model, 
and the Barab\'asi--Albert models.

\section{Graphs: Nodes, Edges, and Adjacency}
\label{sec:graph_basics}

A graph is a set of vertices and a collection of edges,
where each edge joins a pair of distinct vertices. The vertices
correspond to the objects we model, and the edges indicate some
relation between pairs of objects. Throughout this dissertation
I restrict attention to simple, undirected, finite graphs
\cite{hofstad2017, bollobas2001}.

\begin{definition}[Graph]
  A \emph{graph} $G = (V, E)$ consists of a finite vertex set $V$
  with $|V| = n$ and an edge set
  $E \subseteq \bigl\{\{u,v\} : u,v \in V,\; u \neq v\bigr\}$.
\end{definition}

A graph is called \emph{simple} if it contains no self-loops
(edges from a vertex to itself) and no multi-edges (more than
one edge between the same pair of vertices). A graph is called
\emph{undirected} if every edge $\{u,v\}$ is unordered, meaning
that a connection from $u$ to $v$ implies an equally valid
connection from $v$ to $u$. All graphs in this dissertation are
simple and undirected unless stated otherwise.

The \emph{neighbours} of a vertex $v$ are the vertices directly
connected to $v$ by an edge, that is, the set
$N(v) = \{u \in V : \{u,v\} \in E\}$.
A \emph{path} between two vertices $u$ and $v$ is a sequence of
distinct vertices $u = w_0, w_1, \dots, w_\ell = v$ such that
$\{w_{i-1}, w_i\} \in E$ for each $i$; the \emph{length} of the
path is the number of edges $\ell$.

A \emph{random graph} is a probability distribution over the set
of all graphs on a fixed vertex set $V$; equivalently, it is a
graph-valued random variable. Each realisation drawn from this
distribution is an ordinary deterministic graph, but the
ensemble of all realisations encodes a probability distribution
over graph structures \cite{hofstad2017}. The three models
studied in Chapters~4--6 each define such a distribution in a
specific way.

The adjacency structure of a graph can be represented by an 
$n \times n$ matrix. Fix an ordering of the vertices, and 
define the \emph{adjacency matrix} $A$ by
\[
  A_{ij} =
  \begin{cases}
    1, & \text{if } \{i,j\} \in E,\\
    0, & \text{otherwise}.
  \end{cases}
\]
For a simple undirected graph, $A$ is symmetric with zeros on
the diagonal.

The \emph{degree} of a vertex $v$, denoted $\deg(v)$, is the 
number of neighbours of $v$, i.e.\ $\deg(v) = |N(v)|$.
The \emph{degree sequence} of a graph is the list
\[
  \bigl(\deg(v_1), \deg(v_2), \dots, \deg(v_n)\bigr).
\]
The \emph{mean degree} is defined by
\[
  \langle k \rangle
  = \frac{1}{n} \sum_{v \in V} \deg(v)
  = \frac{2|E|}{n},
\]
where the identity $\sum_v \deg(v) = 2|E|$ holds because each
edge contributes one to the degree of each of its two endpoints.

In the Erd\H{o}s--R\'enyi model $G(n,p)$, edges are present 
independently with probability $p \in [0,1]$. For this 
model the expected mean degree satisfies
\[
  \mathbb{E}[\langle k \rangle] = (n-1)p \approx np
\]
when $n$ is large. This quantity plays a central role in 
characterising the connectivity properties and phase 
transition behaviour of random graphs.

A \emph{connected component} of $G$ is a maximal set of
vertices $S \subseteq V$ such that every pair of vertices in $S$
is joined by a path contained entirely within $S$. By
\emph{maximal} we mean that no additional vertex can be added to
$S$ while preserving this property. A graph is called
\emph{connected} if it has exactly one connected component,
i.e.\ a path exists between every pair of vertices.

As a concrete example, consider a graph with four vertices 
$V = \{a,b,c,d\}$ and edges
\[
  E = \bigl\{\{a,b\},\, \{b,c\},\, \{c,d\}\bigr\}.
\]
The adjacency matrix is
\[
  A =
  \begin{pmatrix}
    0 & 1 & 0 & 0 \\
    1 & 0 & 1 & 0 \\
    0 & 1 & 0 & 1 \\
    0 & 0 & 1 & 0
  \end{pmatrix},
\]
the degree sequence is $(1,2,2,1)$, and the mean degree 
is $\langle k \rangle = 6/4 = 1.5$. This graph is connected,
since a path exists between every pair of vertices.

\section{Key Metric Definitions}
\label{sec:metrics}

To compare the three random graph models, we will use 
a set of network metrics that appears in figures in
the results chapters. This section defines the degree 
distribution, clustering coefficient, average path length, 
giant component fraction, and graph density
\cite{hofstad2017, albert2002}.

\paragraph{Degree distribution.}
The \emph{degree distribution} $P(k)$ of a graph is the 
fraction of vertices with degree exactly $k$:
\[
  P(k) = \frac{1}{n} \bigl|\{v \in V : \deg(v) = k\}\bigr|.
\]
This summarises how connectivity is spread across the network.
For large $k$, the behaviour of $P(k)$ is called the
\emph{tail} of the distribution. Networks whose tails decay
slowly --- so that very high-degree vertices are far more
common than in a normal distribution --- are called
\emph{heavy-tailed}, and are associated with the presence of
highly connected hubs. This tail behaviour distinguishes
Poisson-like networks from heavy-tailed, scale-free networks
(scale-free networks are defined precisely in
Section~\ref{sec:probability}).

\paragraph{Clustering coefficient.}
A \emph{clique} is a set of vertices that are all pairwise
connected to one another; the clustering coefficient measures
how close a vertex's neighbourhood is to forming a clique.
The local clustering coefficient $C(v)$ of a vertex 
$v$ with $\deg(v) \ge 2$ is
\[
  C(v) =
  \frac{\text{number of edges among neighbours of } v}
       {\binom{\deg(v)}{2}},
\]
where $\binom{\deg(v)}{2} = \frac{\deg(v)(\deg(v)-1)}{2}$ is
the total number of pairs of neighbours, i.e.\ the maximum
possible number of edges among them.
By convention $C(v) = 0$ if $\deg(v) \le 1$. The 
\emph{average clustering coefficient} of the graph is
\[
  C = \frac{1}{n} \sum_{v \in V} C(v).
\]
This captures the tendency of a network to contain 
tightly connected local neighbourhoods \cite{hofstad2017}.

\paragraph{Average path length.}
For two vertices $u$ and $v$ that are connected by at least one
path, the \emph{distance} $d(u,v)$ is the length of a shortest
path between them. The 
\emph{average path length} $L$ is defined by
\[
  L = \frac{1}{n(n-1)} \sum_{u \ne v} d(u,v),
\]
where the sum is taken over all ordered pairs $(u,v)$
belonging to the largest connected component of the graph. 
This tells you how many steps are required on average
to travel between two randomly chosen vertices
\cite{hofstad2017, albert2002}.

\paragraph{Giant component fraction.}
The \emph{largest connected component} of a graph is the
connected component (defined in Section~\ref{sec:graph_basics})
with the greatest number of vertices. When its size grows
proportionally to $n$ as $n \to \infty$, it is called the
\emph{giant component}. The \emph{giant component fraction} is
\[
  \phi = \frac{|V_{\text{giant}}|}{n},
\]
where $V_{\text{giant}}$ denotes the vertex set of the 
largest connected subgraph. This metric is central to 
studying connectivity phase transitions in the 
Erd\H{o}s--R\'enyi model \cite{hofstad2017}.

\paragraph{Graph density.}
The \emph{density} $\rho$ of an undirected simple graph 
with $n$ vertices and $|E|$ edges is
\[
  \rho = \frac{2|E|}{n(n-1)},
\]
which measures the fraction of possible edges that are 
actually present. The sparser the network, the closer $\rho$ 
is to zero, while denser networks have $\rho$ closer to one
\cite{hofstad2017}.

These metrics provide a means of comparing the 
Erd\H{o}s--R\'enyi, Watts--Strogatz, and Barab\'asi--Albert 
models and for answering the research questions.

\section{The Graph Laplacian and Spectral Connectivity}
\label{sec:laplacian}

The adjacency matrix captures which pairs of vertices are
directly connected, but many global connectivity properties
are better expressed by the graph Laplacian. The Laplacian
links graph structure to linear algebra by encoding degrees
and adjacencies in a single symmetric matrix \cite{hofstad2017}.

Given a simple graph $G = (V,E)$ with adjacency matrix $A$, 
the \emph{degree matrix} $D$ is the diagonal matrix with entries
\[
  D_{ii} = \deg(v_i)
\]
for each vertex $v_i$. The \emph{combinatorial graph Laplacian}
$L$ is defined by
\[
  L = D - A.
\]
The word \emph{combinatorial} here distinguishes this matrix
from other Laplacian variants (such as the normalised
Laplacian); it is the most natural version for counting
arguments. For an undirected graph, $L$ is \emph{symmetric}
(i.e.\ $L = L^T$) and \emph{positive semi-definite}, meaning
that all of its eigenvalues are real and non-negative.
An \emph{eigenvalue} of $L$ is a scalar $\lambda$ for which
there exists a non-zero vector $\mathbf{x}$ (called an
\emph{eigenvector}) satisfying $L\mathbf{x} = \lambda\mathbf{x}$.

Let the eigenvalues of $L$ be ordered as
\[
  0 = \lambda_1 \le \lambda_2 \le \cdots \le \lambda_n.
\]
The smallest eigenvalue $\lambda_1 = 0$ always exists, with 
corresponding eigenvector given by the all-ones vector
$\mathbf{1} = (1, 1, \dots, 1)^T$, since every row of $L$
sums to zero by construction. The 
second-smallest eigenvalue $\lambda_2$ is called the
\emph{algebraic connectivity} or \emph{Fiedler value}, and it
provides a quantitative measure of how well connected the graph
is. A classical result states that $\lambda_2 > 0$ if and only
if the graph is connected, so the Laplacian spectrum encodes
connectivity information in a way that is directly accessible
using linear algebra techniques.

This spectral viewpoint provides the rigorous link between 
the linear algebra modules and the study of random graphs 
in this project. In particular, the Laplacian and its 
eigenvalues appear when interpreting the robustness and 
connectivity properties of networks generated by the three 
random graph models.

\section{Probabilistic Groundwork}
\label{sec:probability}

All three random graph models used in this dissertation are 
probabilistic objects: edges are present or absent according to 
specified random mechanisms. This section briefly recalls the 
basic probability distributions and limit theorems required to 
analyse their degree distributions and connectivity properties
\cite{hofstad2017, bollobas2001}.

A \emph{Bernoulli trial} is a random experiment with exactly two
outcomes, usually coded as $X = 1$ (success) and $X = 0$
(failure), with probabilities
\[
  \mathbb{P}(X = 1) = p
  \quad\text{and}\quad
  \mathbb{P}(X = 0) = 1 - p.
\]
The \emph{binomial distribution} $\mathrm{Bin}(n,p)$ describes 
the total number of successes in $n$ independent Bernoulli
trials with the same success probability $p$. Its
\emph{probability mass function} (the function giving the
probability of each specific outcome) is
\[
  \mathbb{P}(X = k)
  = \binom{n}{k} p^k (1-p)^{n-k},
  \qquad k = 0,1,\dots,n.
\]
In the Erd\H{o}s--R\'enyi model $G(n,p)$, the degree of a fixed 
vertex has distribution $\mathrm{Bin}(n-1,p)$, since each of the 
other $n-1$ vertices is independently joined to it with 
probability $p$.

The \emph{Poisson limit theorem} states that if $n \to \infty$ 
and $p = p_n$ is chosen so that $\lambda = np_n$ remains
fixed, then $\mathrm{Bin}(n,p_n)$ becomes increasingly
well-approximated by a Poisson random variable with parameter
$\lambda$. The Poisson 
distribution has probability mass function
\[
  \mathbb{P}(X = k)
  = e^{-\lambda} \frac{\lambda^k}{k!},
  \qquad k = 0,1,2,\dots.
\]
This result applies to \emph{sparse} Erd\H{o}s--R\'enyi graphs,
meaning graphs where $p = \lambda/n$ for a fixed constant
$\lambda$, so that the mean degree $\lambda = np$ stays bounded
as $n$ grows rather than increasing with the number of vertices.
In this regime the degree of each vertex approximates a
Poisson random variable with mean $\lambda$.

Finally, a \emph{power-law distribution} is one whose tail
behaves like
\[
  \mathbb{P}(X = k) \sim k^{-\tau}
\]
for some exponent $\tau > 1$ as $k \to \infty$. Networks whose 
degree distributions follow a power law are often called
\emph{scale-free}, because the shape of the distribution looks
similar regardless of the scale at which degrees are measured:
multiplying $k$ by any constant changes $P(k)$ only by a
multiplicative factor, preserving the overall form. The
Barab\'asi--Albert model is an example of a random graph model
that produces approximate power-law degree distributions with
exponent $\tau \approx 3$ \cite{albert2002}.

The formal definitions in this chapter follow the treatment
in \cite{hofstad2017} and \cite{bollobas2001}; the metric
definitions in Section~\ref{sec:metrics} are also used
throughout \cite{albert2002}.

% ════════════════════════════════════════════════════════════
%  BIBLIOGRAPHY
% ════════════════════════════════════════════════════════════
\begin{thebibliography}{99}

\bibitem{hofstad2017}
  R.~van der Hofstad,
  \textit{Random Graphs and Complex Networks, Volume~I}.
  Cambridge University Press, 2017.

\bibitem{watts1998}
  D.~J.~Watts and S.~H.~Strogatz,
  ``Collective dynamics of `small-world' networks,''
  \textit{Nature}, \textbf{393}, 440--442, 1998.

\bibitem{barabasi1999}
  A.-L.~Barab\'{a}si and R.~Albert,
  ``Emergence of scaling in random networks,''
  \textit{Science}, \textbf{286}(5439), 509--512, 1999.

\bibitem{erdos1959}
  P.~Erd\H{o}s and A.~R\'{e}nyi,
  ``On random graphs~I,''
  \textit{Publicationes Mathematicae Debrecen}, \textbf{6},
  290--297, 1959.

\bibitem{networkx}
  A.~A.~Hagberg, D.~A.~Schult, and P.~J.~Swart,
  ``Exploring network structure, dynamics, and function using
  NetworkX,''
  in \textit{Proc.\ 7th Python in Science Conference (SciPy 2008)},
  pp.~11--15, 2008.

\bibitem{albert2002}
  R.~Albert and A.-L.~Barab\'{a}si,
  ``Statistical mechanics of complex networks,''
  \textit{Reviews of Modern Physics}, \textbf{74}, 47--97, 2002.

\bibitem{bollobas2001}
  B.~Bollob\'{a}s,
  \textit{Random Graphs}, 2nd edn.
  Cambridge University Press, 2001.

\end{thebibliography}

\end{document}
